% Results section (draft, 2026-10-14). Needs in the preamble: booktabs, \CropCode, \todo, and bib keys
% oquab2023dinov2 / bachmann2022multimae (the other keys are already cited in method.tex).
% Numbers are read from reports/*.csv and outputs/*/training_history.json;
% the file that produced each table is named in a comment above it.
%
% DATA VERSIONS. Every sorghum ablation below (E2, controls, E3, E4, E8, partial clouds, E9)
% ran on the first release of the sorghum set (grey-background RGB, waviness leaf fields) and
% is being re-run on the final one (background-free, brightened RGB; per-leaf width); the
% four re-run arms finished so far reproduce the modality ordering (Chamfer 5.00 / 1.85 / 2.53
% x 1e-3 for PC / PC+RGB+D / PC+RGB+D+params). Maize E2 and E3 are on the final maize set;
% maize E4, E8 and the pose rows are on its first release.

\section{Experiments}
\label{sec:results}

\paragraph{Setup.}
All ablations share one budget of 197{,}400 optimizer steps at global batch 32
(600 epochs of one view per plant over the 10{,}500 training plants), one
view-sampling seed, and the base ViT (114\,M parameters) unless the model size
is the variable; arms that change the number of training plants keep the step
budget by changing their epoch count (Sec.~\ref{sec:pretrain}). The headline
sorghum model (E1) differs only in schedule: 1{,}000 epochs over all ten views
at global batch 256. Phenotype readout follows Sec.~\ref{sec:field}: ridge
regression on the frozen CLS token, the parameter tensor zeroed, one
deterministic view per plant, reported as $R^2$ on the extreme-enriched
held-out plants (2{,}250 val, 2{,}250 test; val unless stated) with paired
bootstrap confidence intervals over plants (2{,}000 resamples). A difference
counts as a \emph{win} only when its interval excludes zero on val \emph{and}
test. Reconstruction is reported as symmetric Chamfer distance in the
unit-sphere frame ($\times10^{-3}$; a plant is $\approx$1.4\,m, so 1.0 is
about 7\,mm) on the masked validation step. Sorghum's four named targets have
rank two by construction (height and leaf count are one generator variable,
$r{=}0.994$, and the biomass proxy is that variable again), so we add the
independent axes the generator actually varies, leaf roll and leaf length; maize
exposes eleven largely independent traits (Sec.~\ref{sec:data}).

\subsection{Which streams carry the phenotype}
\label{sec:res-e2}

% reports/probe_pr_600_16592054.csv (val/test), outputs/e2_*/training_history.json,
% reports/probe_pm_600_22751615.csv (maize, final data)
\begin{table}[t]
  \centering\scriptsize
  \setlength{\tabcolsep}{3.2pt}
  \begin{tabular}{l c cccccc}
    \toprule
    & Chamfer & \multicolumn{6}{c}{frozen-probe $R^2$ (val)} \\
    \cmidrule(lr){3-8}
    streams & $\times10^{-3}$ & height & leaves & biomass & branch & roll & leaf len \\
    \midrule
    \multicolumn{8}{l}{\emph{Sorghum}} \\
    PC                 & 4.96 & 0.899 & 0.902 & 0.900 & 0.814 & 0.452 & 0.102 \\
    PC+RGB             & 2.50 & 0.978 & 0.983 & 0.985 & 0.918 & 0.796 & 0.273 \\
    PC+RGB+D           & \textbf{1.71} & \textbf{0.981} & \textbf{0.984} & \textbf{0.986} & 0.913 & \textbf{0.806} & \textbf{0.362} \\
    PC+RGB+D+$\theta$  & 2.50 & 0.966 & 0.967 & 0.969 & 0.899 & 0.793 & 0.280 \\
    \midrule
    & & height & leaves & biomass & angle & curl & stem r. \\
    \multicolumn{8}{l}{\emph{Maize}} \\
    PC                 & 2.63 & 0.171 & 0.670 & 0.610 & 0.756 & 0.114 & 0.289 \\
    PC+RGB             & 1.61 & 0.620 & 0.914 & \textbf{0.783} & \textbf{0.927} & \textbf{0.334} & 0.753 \\
    PC+RGB+D           & 1.21 & 0.633 & 0.914 & 0.774 & 0.926 & 0.236 & \textbf{0.760} \\
    PC+RGB+D+$\theta$  & \textbf{0.99} & \textbf{0.635} & 0.908 & 0.781 & 0.919 & 0.304 & 0.743 \\
    \bottomrule
  \end{tabular}
  \caption{\textbf{Modality value-add (E2).} Four arms per species at equal
  steps, differing only in which token streams exist. $\theta$ is the
  parameter stream as an \emph{input} modality inside the shared masking
  budget. Probe $R^2$ on held-out plants; branch and roll are the two sorghum
  leaf-angle definitions, leaf len is the per-plant mean leaf length. Test
  values are within 0.03 of val on every cell.}
  \label{tab:e2}
\end{table}

Table~\ref{tab:e2} isolates each stream at equal compute. Three facts hold in
both species.

\textbf{RGB is where the phenotype arrives.} Adding the photograph to the point
cloud is the largest move in every column: sorghum height $0.899\!\to\!0.978$,
roll $0.45\!\to\!0.80$; maize height $0.17\!\to\!0.62$, stem radius
$0.29\!\to\!0.75$. A cloud alone carries the silhouette but not the rendered
cues (shading, leaf overlap, texture) that place organs on the stem.

\textbf{Depth is a reconstruction win and a phenotype wash.} Depth lowers the
sorghum Chamfer from $2.50$ to $1.71$ ($-32\,\%$) and maize from $1.61$ to
$1.21$, but moves the sorghum probe by at most $+0.01$ on size targets and
$+0.09$ on leaf length; on maize the paired comparison gives depth two wins
(stem radius, tassel droop) and four losses (leaf count, droop, twist, width),
none larger than $0.07$. Depth tells the decoder where surfaces are; it does
not tell the encoder more about the plant.

\textbf{The parameter stream, used as an input, does not help.} On sorghum
the four-stream arm loses to PC+RGB+D on every learnable target, on val and
test alike: size targets $-0.015$ to $-0.017$, roll $-0.013$, leaf length
$-0.082$ (all six intervals exclude zero), and its Chamfer gives back exactly
what depth gained ($2.50$ vs $1.71$). On the first maize release the same arm
lost 10 of 11 targets to PC+RGB+D; on the final release it is a tie (height
$0.635$ vs $0.633$, leaf angle $0.919$ vs $0.926$, stem radius $0.743$ vs
$0.760$) while winning the Chamfer ($0.99$ vs $1.21$). Either way the stream
that encodes the answer verbatim at training time is not the stream that
makes the frozen feature better at deployment, where it is absent.

\subsection{The parameter stream: input, target, or neither}
\label{sec:res-control}

% reports/probe_pr_tg100_600_16686383.csv, reports/control_delta_16686383.csv,
% reports/control_delta_vs_tg_16686383.csv
\begin{table}[t]
  \centering\scriptsize
  \setlength{\tabcolsep}{3.5pt}
  \begin{tabular}{l cc c cccc}
    \toprule
    & vision & $\theta$ & Chamfer & \multicolumn{4}{c}{$R^2$ (val)} \\
    \cmidrule(lr){5-8}
    arm & tokens & visible & $\times10^{-3}$ & height & leaves & roll & leaf len \\
    \midrule
    PC+RGB+D                      & 117.0 & --      & 1.71 & \textbf{0.981} & \textbf{0.984} & \textbf{0.806} & \textbf{0.362} \\
    $+\theta$, shared budget      & 107.4 & 58\,\%  & 2.50 & 0.966 & 0.967 & 0.793 & 0.280 \\
    $+\theta$, gated 20\,\%       & 117.0 & 20\,\%  & 2.41 & 0.958 & 0.960 & 0.764 & 0.288 \\
    $+\theta$, gated 60\,\%       & 117.0 & 60\,\%  & 2.72 & 0.922 & 0.925 & 0.676 & 0.202 \\
    $+\theta$, target only        & 117.0 & 0\,\%   & 2.54 & 0.970 & 0.973 & 0.737 & 0.259 \\
    \bottomrule
  \end{tabular}
  \caption{\textbf{Token-budget control (sorghum).} The four-stream arm
  re-trained with the parameter tokens taken out of the shared Dirichlet
  budget, so the three sensor streams see exactly the visible-token count of
  the three-stream arm (117.0, seed-matched masks identical in 5{,}000 of
  5{,}000 draws), at three parameter visibilities. ``Target only'' never
  encodes a parameter token.}
  \label{tab:control}
\end{table}

A fourth stream at mask ratio $0.8$ splits the visible budget four ways, so
Tab.~\ref{tab:e2} cannot tell a token-budget handicap from a representational
one: as trained, the four-stream arm's sensors see $107.4$ visible tokens
against $117.0$, and its parameter tokens are 58\,\% visible. Table~\ref{tab:control}
removes the handicap. Giving the sensor streams back their full budget does
not recover the three-stream probe at any parameter visibility; every gated
arm loses all seven learnable targets to PC+RGB+D on both splits. The trend
with visibility is clean on the size targets (60\,\% $<$ 20\,\% $<$ 0\,\% visible),
and making the parameters a pure reconstruction target still costs $-0.011$ on
size, $-0.07$ on roll and $-0.10$ on leaf length. The loss on roll and leaf
length, the two traits the leaf tokens spell out per leaf, is the largest in
the table, which is consistent with the encoder learning to route those traits
through tokens that will not exist in the field. This is the evidence behind
the design in Sec.~\ref{sec:param-stream}: the recipe is decoded from the
sensor streams and never encoded.

\subsection{Scaling data and capacity}
\label{sec:res-scaling}

% reports/probe_pr_{6168,2024,624,600}_*.csv, outputs/e3_*,e4_*/training_history.json;
% maize: reports/probe_pm_{6168,2024,624}_2275161*.csv (E3, final data), probe_pm_600L_16592061.csv (E4, first release)
\begin{table}[t]
  \centering\scriptsize
  \setlength{\tabcolsep}{3.0pt}
  \begin{tabular}{l r c ccc c c ccc}
    \toprule
    & & \multicolumn{4}{c}{Sorghum} & & \multicolumn{4}{c}{Maize} \\
    \cmidrule(lr){3-6}\cmidrule(lr){8-11}
    & & Chamfer & height & roll & leaf len & & Chamfer & height & angle & stem r. \\
    \midrule
    \multicolumn{11}{l}{\emph{Training plants (E3), base model}} \\
    1{,}000  & 6{,}169 ep & 3.02 & 0.963 & 0.770 & 0.269 & & 3.54 & 0.393 & 0.851 & 0.593 \\
    3{,}000  & 2{,}100 ep & 2.61 & 0.963 & 0.764 & 0.255 & & 1.46 & 0.524 & 0.867 & 0.624 \\
    10{,}000 & 631 ep     & 2.42 & 0.962 & 0.789 & 0.301 & & 0.87 & 0.427 & 0.853 & 0.717 \\
    10{,}500 & 600 ep     & 2.50 & 0.966 & 0.793 & 0.280 & & 0.99 & 0.635 & 0.919 & 0.743 \\
    \midrule
    \multicolumn{11}{l}{\emph{Model size (E4), all plants}} \\
    small  & 25\,M  & 3.02 & 0.946 & 0.761 & 0.252 & & 1.39 & 0.593 & 0.878 & 0.755 \\
    base   & 114\,M & 2.50 & 0.966 & 0.793 & 0.280 & & 1.26 & 0.579 & 0.865 & 0.741 \\
    large  & 332\,M & 2.53 & 0.967 & 0.799 & 0.298 & & 1.19 & 0.629 & 0.893 & 0.740 \\
    \bottomrule
  \end{tabular}
  \caption{\textbf{Data and model scaling at equal steps.} Every row is
  197{,}400 steps; epochs differ so that steps do not. Chamfer $\times10^{-3}$
  and probe $R^2$ (val) of the four-stream arm. The 10{,}000 and 10{,}500-plant
  rows differ by 5\,\% in data and are the run-to-run error bar. Maize E3 is on
  the final maize release; maize E4 on the first (its base row is therefore not
  the E3 full-data row).}
  \label{tab:scaling}
\end{table}

Two near-duplicate runs, 10{,}000 and 10{,}500 training plants, differ by
3\,\% in Chamfer and by $0.004$ in probe $R^2$; we take that as the noise
floor. Against it (Tab.~\ref{tab:scaling}), sorghum reconstruction improves
from 1{,}000 to 3{,}000 plants ($-14\,\%$, five times the floor) and more
weakly to 10{,}000 ($-7\,\%$), while the probe is flat: the size targets sit at
$0.96$ from 1{,}000 plants up, and only leaf length ($0.27\!\to\!0.30$) moves.
Capacity buys more than data at this budget: small to base is $-17\,\%$ in
Chamfer and $+0.02$ to $+0.03$ on every probe target, and large adds nothing
to Chamfer ($2.53$ vs $2.50$) and $+0.01$ to the probe. Maize is noisier. Its
Chamfer falls fourfold from 1{,}000 to 10{,}000 plants, but the probe is not
monotone in plants (height $0.39 / 0.52 / 0.43 / 0.64$) and the three model
sizes are within $0.05$ of each other on every trait; a single seed per arm
does not resolve it. \todo{Paired CIs over the maize E3 arms; a second seed of
\texttt{e3\_10k}.}

\subsection{Baselines}
\label{sec:res-e8}

% reports/probe_e8_base_{sorghum,maize}_*.csv, probe_pr_mean600_16614173.csv, probe_pm_mean600_16614174.csv,
% probe_e8_score_16676974_*.csv, e8_supervised_delta_16676974.csv. Mean pooling for every model.
\begin{table}[t]
  \centering\scriptsize
  \setlength{\tabcolsep}{2.6pt}
  \begin{tabular}{l ccc cc c cccc}
    \toprule
    & \multicolumn{5}{c}{Sorghum} & & \multicolumn{4}{c}{Maize} \\
    \cmidrule(lr){2-6}\cmidrule(lr){8-11}
    encoder (frozen, mean-pooled) & height & leaves & branch & roll & leaf len & & height & angle & curl & stem r. \\
    \midrule
    DINOv2 ViT-B/14 (RGB)~\cite{oquab2023dinov2}      & 0.917 & 0.922 & 0.850 & 0.563 & 0.224 & & 0.547 & 0.812 & 0.083 & 0.518 \\
    MultiMAE-B (RGB+D)~\cite{bachmann2022multimae}     & 0.926 & 0.930 & 0.823 & 0.673 & 0.223 & & 0.535 & 0.772 & 0.069 & 0.439 \\
    Point-MAE (PC)~\cite{pang2022pointmae}             & 0.958 & 0.963 & 0.864 & 0.726 & 0.176 & & 0.520 & 0.870 & 0.189 & 0.438 \\
    EmbodiedMAE-B (RGB+D+PC)~\cite{dong2025embodiedmae}& 0.977 & 0.980 & 0.877 & 0.764 & 0.287 & & 0.517 & 0.822 & 0.169 & 0.415 \\
    supervised from scratch (ours, 4 streams)          & 0.972 & 0.977 & 0.919 & 0.739 & 0.273 & & -- & -- & -- & -- \\
    \midrule
    \CropCode{} PC only                                & 0.910 & 0.914 & 0.820 & 0.512 & 0.106 & & 0.184 & 0.810 & 0.170 & 0.541 \\
    \CropCode{} PC+RGB+D                               & \textbf{0.982} & \textbf{0.985} & 0.912 & \textbf{0.843} & \textbf{0.391} & & \textbf{0.604} & \textbf{0.911} & 0.177 & 0.693 \\
    \CropCode{} four streams                           & 0.971 & 0.973 & 0.902 & 0.810 & 0.330 & & 0.597 & 0.880 & 0.150 & \textbf{0.757} \\
    \bottomrule
  \end{tabular}
  \caption{\textbf{Frozen-encoder baselines (E8).} Published checkpoints read
  the same held-out plants through the same ridge probe; mean pooling for
  every model, the one pooling all share. The supervised row trains our
  architecture end to end on the probe targets and is read through the same
  frozen probe. Maize baselines and the maize \CropCode{} rows are on the
  first maize release; the maize PC-only and PC+RGB+D rows use the CLS token.
  \todo{maize mean-pooled rows for PC / PC+RGB+D, supervised row on maize.}}
  \label{tab:e8}
\end{table}

Table~\ref{tab:e8} reads published encoders through the same probe. On sorghum
the three-stream \CropCode{} wins every learnable target against every
baseline. Against the strongest, EmbodiedMAE-B, the paired margins are
$+0.005$ on the size targets (tight intervals at a $0.98$ ceiling), $+0.08$ on
roll and $+0.10$ on leaf length. Two of the losses are instructive. Point-MAE
beats our PC-only arm on every sorghum target (height $0.958$ vs $0.910$): our
margin comes from fusing a photograph with geometry, not from a better
point-cloud encoder. And the four-stream arm loses the size targets to
EmbodiedMAE-B by $0.005$ to $0.007$, the same parameter-stream penalty as in
Tab.~\ref{tab:e2}, now measured against an external model.

\paragraph{Pretraining against supervision.}
A network of our architecture trained from scratch on the probe targets
reaches $R^2$ $0.974$ height, $0.979$ leaf count, $0.929$ branch angle and
$0.811$ roll through its own head (val), and its CLS probe matches that head
within $0.008$. The frozen three-stream feature beats it on six of seven
targets (size $+0.006$, leaf length $+0.08$ to $+0.12$ on both splits) and
ties on branch angle and roll; the frozen four-stream feature loses size
($-0.008$ to $-0.014$) and the angles ($-0.03$). Pretraining pays without the
parameter stream, and with it the frozen feature is slightly worse than no
pretraining.

\paragraph{Maize and camera pose.}
On the first maize release, \CropCode{} beats EmbodiedMAE-B on 9 of 11
traits and the best baseline per trait on 5 of 11, losing leaf count and curl
to Point-MAE and tassel droop to DINOv2 (the one trait where every
EmbodiedMAE-design encoder is below zero). Point-MAE given each view's camera
pose, so that clouds are gravity-aligned and the plant's azimuth is canonical,
beats our camera-frame four-stream arm on 8 of 11 (height $0.720$ vs
$0.597$). Most of that is a generator artefact: the maize renderer places every
plant's leaf plane in the same world plane, so a full pose rotation hands over
an orientation no rig provides. With gravity only (the camera levelled, its
azimuth kept), Point-MAE's height drops to $0.588$ and camera-frame PC+RGB+D
wins five traits (height, biomass, twist, width, stem radius $+0.23$) and
loses four (leaf angle, leaf count, curl $-0.18$, tassel droop $-0.21$). Our
frozen fused arms gain nothing from levelled clouds (0 wins, 5 losses against
their own camera-frame rows); PC-only gains five. \todo{Arm pretrained on
levelled clouds (\texttt{maize\_e2\_pcrgbd\_levelled}) was cancelled at epoch 9;
decide whether to report it.}

\subsection{What a camera actually sees}
\label{sec:res-views}

% reports/pcview_16676994_{coverage,recon,probe,bootstrap}.csv (sorghum, first release)
\begin{table}[t]
  \centering\scriptsize
  \setlength{\tabcolsep}{3.4pt}
  \begin{tabular}{l ccc ccc ccc}
    \toprule
    & \multicolumn{3}{c}{matched probe, height} & \multicolumn{3}{c}{transferred probe, height} & \multicolumn{3}{c}{Chamfer $\times10^{-3}$} \\
    \cmidrule(lr){2-4}\cmidrule(lr){5-7}\cmidrule(lr){8-10}
    cloud & full & 3 cam & 1 cam & full & 3 cam & 1 cam & full & 3 cam & 1 cam \\
    coverage & 100\,\% & 78\,\% & 49\,\% & & & & & & \\
    \midrule
    PC             & 0.899 & 0.857 & 0.807 & 0.899 & 0.615 & 0.633 & 4.94 & 6.16 & 13.3 \\
    PC+RGB         & 0.978 & 0.968 & 0.951 & 0.978 & 0.690 & $-0.06$ & 2.42 & 3.59 & 11.3 \\
    PC+RGB+D       & 0.980 & 0.973 & 0.963 & 0.980 & 0.764 & 0.316 & 1.65 & 2.97 & 11.0 \\
    four streams   & 0.965 & 0.956 & 0.948 & 0.965 & 0.779 & 0.753 & 2.43 & 3.67 & 11.2 \\
    E1 (1k ep)     & 0.983 & 0.977 & 0.964 & 0.983 & 0.670 & $-0.27$ & 1.31 & 2.74 & 10.2 \\
    \bottomrule
  \end{tabular}
  \caption{\textbf{Partial clouds (sorghum, test).} Every model trained on
  complete clouds; at evaluation the cloud is cut to the points visible from
  camera 0, or from cameras 0--2, by ray-casting against the mesh. Matched: the
  probe refit per condition. Transferred: the full-cloud probe applied as is.
  Chamfer against the complete cloud, masked regime, val.}
  \label{tab:views}
\end{table}

Training clouds are complete; a rig's are not. Table~\ref{tab:views} scores
each model on the subset of points one or three cameras actually see. The
information survives: refit per condition, a model that sees RGB loses only
$0.01$ to $0.03$ on size at one camera and under $0.01$ at three (PC-only
$-0.09$ and $-0.04$), and roll and leaf length lose $0.02$ to $0.06$. Depth
matters more on a partial cloud than on a full one: at one camera PC+RGB loses
$-0.027\ [-0.030, -0.024]$ on height and PC+RGB+D $-0.017\ [-0.019, -0.015]$.
But the latent moves with coverage. A readout fitted on complete clouds and
transferred to one-camera clouds collapses (E1 height $0.983 \to 0.670 \to -0.27$),
and refitting recovers $0.964$. Any probe calibrated on complete scans is
biased on real ones, which is why the real-data protocol (Sec.~\ref{sec:field})
reports the transfer gap. The decoder does not complete either: at one camera
the distance from the true cloud to the prediction is $8.7$ to $10.4$ against
$4.0$ for echoing the partial input, so the model reproduces what it is given
rather than hallucinating the hidden half.

\subsection{Single-sensor generation and distillation}
\label{sec:res-distill}

% outputs/4m_distill_15k_all/{warmstart_eval,training_history}.json, reports/quant_*_{val,test}.json,
% outputs/maize_distill_all (first maize release; run stopped at epoch 59 of 100), reports/single_source_pc_*.csv
\begin{table}[t]
  \centering\scriptsize
  \setlength{\tabcolsep}{4pt}
  \begin{tabular}{l cc cc c cc cc}
    \toprule
    & \multicolumn{4}{c}{Sorghum (100 ep)} & & \multicolumn{4}{c}{Maize (55 ep)} \\
    \cmidrule(lr){2-5}\cmidrule(lr){7-10}
    & \multicolumn{2}{c}{Chamfer $\times10^{-3}$} & \multicolumn{2}{c}{$\theta$ MAE} & & \multicolumn{2}{c}{Chamfer $\times10^{-3}$} & \multicolumn{2}{c}{$\theta$ MAE} \\
    source & before & after & before & after & & before & after & before & after \\
    \midrule
    RGB         & 0.97 & 0.88 & 0.032 & 0.030 & & 1.80 & 0.98 & 0.078 & 0.074 \\
    depth       & 0.73 & 0.62 & 0.030 & 0.029 & & 1.81 & 0.99 & 0.079 & 0.074 \\
    point cloud & 1.92 & 1.22 & 0.039 & 0.035 & & 51.7 & 1.97 & 0.090 & 0.078 \\
    $\theta$    & 6.87 & 6.34 & --    & --    & & 32.6 & 11.8 & --    & --    \\
    \bottomrule
  \end{tabular}
  \caption{\textbf{Single-source generation before and after distillation.}
  Each row: that stream alone, the other three absent, generating the rest.
  Before = the pretrained model (sorghum E1 at 1{,}000 epochs; maize four-stream
  arm at 600) used zero-shot; after = the distilled student. Parameter MAE on
  the normalized $[0,1]$ fields of hidden real leaf tokens. The maize
  before-row for the point cloud is inflated by the collapse of an unmasked
  cloud (Sec.~\ref{sec:distill}).}
  \label{tab:distill}
\end{table}

The pretrained model already generates cross-modally without any
single-sensor training: on sorghum a photograph alone yields the cloud at
Chamfer $0.97$ and the recipe at MAE $0.032$ (Tab.~\ref{tab:distill}), and the
same holds on maize ($1.80$, $0.078$). Distillation against the full-modality
teacher improves every source, most where the source was weakest: sorghum
$-9\,\%$ / $-15\,\%$ / $-37\,\%$ in Chamfer from RGB / depth / cloud, maize
$-45\,\%$ / $-45\,\%$ / $-96\,\%$, with parameter error down $4$ to $13\,\%$.
Two seeds of the RGB-only student agree to the fourth digit: on the 2{,}250
test plants the pretrained model generates the cloud from a photograph at
$1.00 \times 10^{-3}$ and both students at $0.79$ (val $1.01 \to 0.80$); from a
depth frame the student reaches $0.62$. Expressed as skill against the
mean plant of the split, the photograph-to-cloud student is at $0.984$.
Hiding half the source's tokens during distillation (the regime of
Sec.~\ref{sec:distill}) beats full visibility, $0.80$ vs $0.85$, so the
regularization costs nothing. The maize student's point-cloud row is the
collapse the teacher protocol was designed around: an unmasked cloud fed to the
pretrained model reconstructs at $51.7$, the distilled student at $1.97$.
\todo{Maize distillation on the final release is queued (teacher
\texttt{maize\_4m\_d2}); replace the 55-epoch numbers.} The specialist
comparison on sorghum holds here too: a student distilled on the parameters
alone recovers them from a cloud at MAE $0.039$, the four-target generalist at
$0.035$.

\subsection{What the latent encodes}
\label{sec:res-e9}

% reports/e9_pr_600_16592054/e9_latent_summary_val.csv, reports/e9_pr_tg100_600_16686383/, reports/e9_pm_600_22751615/
There is no genotype in either generator: every plant is one draw from one
continuous distribution, so a clustered latent would be a defect, not a
finding. The frozen CLS space is a continuum (best silhouette $0.17$ to $0.23$
at $k{=}2$ with adjusted Rand index $\approx 0$ against shape deciles, t-SNE
trustworthiness $0.95$ to $0.98$) of low effective dimension (participation
ratio $4.4$ to $6.3$ of 768 on sorghum, $6.7$ to $10.3$ on maize; 14 to 19
principal components for 90\,\% of the variance once RGB is present). Its
leading direction is the plant's size: the first principal component correlates
with leaf count at $|r| = 0.73$ to $0.94$ in the sorghum arms. The parameter
stream shapes this geometry in the direction the probe already showed: the
gated arms are the least size-aligned ($|r|$ $0.48$ and $0.34$) and the most
isotropic (participation ratio $8.5$ and $14.3$), and the target-only arm is
the most size-dominated of the five (first component 43\,\% of variance,
$|r| = 0.95$). \todo{Figure: t-SNE of the three-stream and four-stream latents
colored by leaf count and by roll.}

\subsection{Field data}
\label{sec:res-field}

\todo{E10. Real-scan evaluation of the RGB-only and depth-only students
against the staged-pipeline ground truth (Sec.~\ref{sec:field}): parameter
error, re-rendered geometry, and the one-camera transfer gap of
Tab.~\ref{tab:views} as the predicted sim-to-real penalty. Not yet run.}
